\begin{lemma}[Support of the copied left values]
Let $Q$ be a type, let $r$ be a binary relation on $Q$, and let
$h:\mathsf{IncSeq}\to\mathsf{PowerQ}(Q)$ be a multi-sequence bad for
$\mathsf{PowerRel}(r)$, meaning that for every increasing enumeration $z$,
\[
\neg\mathsf{PowerRel}(r)\bigl(h(z),h(\mathsf{ShiftMap}(z))\bigr).
\]
Then for every increasing enumeration $x$ and all $n,k\in\mathbb N$, the
selected left value in row $n$ has support contained in the support of the
row's initial presentation:
\[
\mathsf{PowerSupport}(\mathsf{PowerCopiedLeft}_r(h,x,n,k))
\subseteq
\mathsf{PowerSupport}\bigl(h(\mathsf{PowerShiftIter}(n,x))\bigr).
\]
\end{lemma}
\begin{proof}
Fix an increasing enumeration $x$ and a row $n$, and write
\[
L_{n,k}=\mathsf{PowerCopiedLeft}_r(h,x,n,k).
\]
We prove the displayed inclusion by induction on $k$.  At depth zero,
set $y=\mathsf{PowerShiftIter}(n,x)$.  The definition
\noderef{BadMultiSequence}, instantiated at this increasing enumeration $y$,
gives the raw failure
\[
\neg\mathsf{PowerRel}(r)\bigl(h(y),h(\mathsf{ShiftMap}(y))\bigr).
\]
The successor equation in \noderef{PowerShiftIter} gives
\[
\mathsf{PowerShiftIter}(n+1,x)=\mathsf{ShiftMap}(\mathsf{PowerShiftIter}(n,x))
  =\mathsf{ShiftMap}(y).
\]
Consequently,
\[
\neg\mathsf{PowerRel}(r)\bigl(
 h(\mathsf{PowerShiftIter}(n,x)),
 h(\mathsf{PowerShiftIter}(n+1,x))\bigr).
\]
The depth-zero clause of the definition \noderef{PowerCopiedLeft} gives
\[
L_{n,0}=\mathsf{PowerIResponse}_r\bigl(
  h(\mathsf{PowerShiftIter}(n,x)),
  h(\mathsf{PowerShiftIter}(n+1,x))\bigr).
\]
Applying \noderef{PowerIResponseLegal} to the preceding failed pair and
using this defining equation gives
\[
\mathsf{PowerMove}\bigl(h(\mathsf{PowerShiftIter}(n,x)),L_{n,0}\bigr).
\]
Now \noderef{PowerMoveInSupport} yields
\[
\mathsf{PowerSupport}(L_{n,0})
\subseteq
\mathsf{PowerSupport}\bigl(h(\mathsf{PowerShiftIter}(n,x))\bigr),
\]
which is the induction base.

For the successor step, assume
\[
\mathsf{PowerSupport}(L_{n,k})
\subseteq
\mathsf{PowerSupport}\bigl(h(\mathsf{PowerShiftIter}(n,x))\bigr).
\]
The second conjunct of \noderef{PowerCopiedFailureMove}, applied at depth
$k$ in row $n$, gives the legal move
\[
\mathsf{PowerMove}(L_{n,k},L_{n,k+1}).
\]
Applying \noderef{PowerMoveInSupport} to this move gives
\[
\mathsf{PowerSupport}(L_{n,k+1})
\subseteq \mathsf{PowerSupport}(L_{n,k}).
\]
For any $q\in\mathsf{PowerSupport}(L_{n,k+1})$, the last inclusion puts $q$
in $\mathsf{PowerSupport}(L_{n,k})$, and the induction hypothesis then puts
$q$ in
$\mathsf{PowerSupport}(h(\mathsf{PowerShiftIter}(n,x)))$.  This is exactly
the desired inclusion at depth $k+1$.  Since $n$ and $x$ were arbitrary, the
result holds for every row and depth.
\end{proof}
